\documentclass[10pt,a4paper]{amsart}
\usepackage[margin=19mm]{geometry}
\usepackage{amsmath,amssymb,mathtools}
\usepackage[scr=rsfs,cal=euler]{mathalfa}
\usepackage{xfrac}
\usepackage[unicode,hidelinks,bookmarks=false]{hyperref}
\newcommand{\ie}{i.e.\ }
\newcommand{\cf}{cf.\ }
\newcommand{\resp}{respectively\ }
\newcommand{\Subsection}{\subsection}
\providecommand{\nobreakdash}{\nobreak\hbox{-}}
\setlength{\emergencystretch}{2em}
\multlinegap0pt
\usepackage{amssymb}
\usepackage{mathtools}
\usepackage{dsfont}
\usepackage{cancel}
\usepackage{xcolor}
\usepackage{enumerate}
\usepackage[shortlabels]{enumitem}
\usepackage{tikz}
\newtheorem{lemma}{Lemma}
\newcommand{\dd}{\,\mathrm{d}}
\newcommand{\defeq}{\coloneqq}
\newcommand{\disp}{\displaystyle}
\newcommand{\e}{\varepsilon}
\renewcommand{\geq}{\geqslant}
\renewcommand{\leq}{\leqslant}
\title{The liquid drop model with a Yukawa potential: existence of minimizers and sharp stability of the ball}
\author{Lia Bronsard, Kenneth DeMason,, Ihsan Topaloglu}
\date{Source: arXiv:2608.19340v1 (CC BY 4.0)}
\begin{document}
\maketitle

\noindent\emph{Source and licence.} The liquid drop model with a Yukawa potential: existence of minimizers and sharp stability of the ball. Version arXiv:2608.19340v1 (CC BY 4.0), distributed by arXiv under the Creative Commons Attribution 4.0 International licence, http://creativecommons.org/licenses/by/4.0/, as recorded in the arXiv licence metadata for this exact version and shown on the abstract page https://arxiv.org/abs/2608.19340v1. Full licence notice: \texttt{licenses/arxiv-2608.19340-CC-BY-4.0.txt}.\par\smallskip

\noindent\textit{Editorial scope.} Counted unit: the article's lemma lem:characterization\_G (source0.tex lines 1617--1626). The article's complete original proof of it is delivered with it (source0.tex lines 1628--1664, one \texttt{proof} environment of 37 lines). No mathematics is written, repaired or re-derived: the statement and the proof are the article's own lines, in the article's own notation, and no retained line is substituted or re-flowed.  Retained uncounted as the source-local context the proof needs: lines 1613--1615.  Nothing was omitted from any retained span, so the package carries no omission record.  No retained line contains a citation, so the wrapper carries no bibliography and no bibliography entry.  No retained line contains a figure environment, an image-inclusion command or drawing code, so no image is delivered and the package declares no asset dependency.  Editorial formatting only, in addition: an AMS \texttt{amsart} preamble that restates the article's own declarations and macros used by the retained lines, namely the environments \texttt{lemma} on separate counters, together with the standard packages those lines need.  The retained environments print in this document as Lemma 1 in order of appearance, whereas the article prints the same items with the numbering of its own full document.  The complete native source of the article as distributed in the arXiv e-print for this exact version is bundled unchanged as \texttt{sources/208077/source0.tex}.
	\begin{equation}\label{eqn:G}
		G(\tau) \defeq \inf \Big\{A(\varphi)+\tau B(\varphi) \colon \varphi\in\mathscr{A}, \ \varphi(1)=1\Big\}
	\end{equation}

	\begin{lemma}\label{lem:characterization_G}
		For any $l\geq 1$, $\disp G(\lambda_l) = \frac{1}{\alpha \, i_l(\alpha) k_l(\alpha)}$,
		and the infimum in \eqref{eqn:G} at $\tau=\lambda_l$ is attained at
			\[
				\varphi_l(r) \defeq \begin{dcases*}
						 				\frac{i_l(\alpha r)}{i_l(\alpha)} &  if $0<r\leq 1$, \\
						 				\frac{k_l(\alpha r)}{k_l(\alpha)}	&  if $r \geq 1$.
									\end{dcases*}
			\]
	\end{lemma}

\begin{proof}
Using the behavior of $i_l$ at $r=0$ and $k_l$ at $r=\infty$ we see that $\varphi_l \in \mathscr{A}$. Also, $\varphi_l$ is continuous at $r=1$ and $\varphi_l(1)=1$. Since $i_l$ and $k_l$ solve the modified spherical Bessel equation $x^2u'' +2xu'-\big(x^2+l(l+1)\big)u=0$, $\varphi_l$ solves $\big(r^2\varphi'\big)'=\big(\lambda_l+\alpha^2r^2\big)\varphi$ on $(0,1)$ and $(1,\infty)$, which is exactly the Euler--Lagrange equation of $Q(\varphi)\defeq A(\varphi)+\lambda_l B(\varphi)$.

Now let $\varphi\in\mathscr{A}$ be arbitrary with $\varphi(1)=1$. Define $w\defeq\varphi-\varphi_l$. Then $w\in\mathscr{A}$ with $w(1)=0$, and $Q(\varphi)=Q(\varphi_l)+2I(\varphi_l,w)+Q(w)$, where
	\[
		I(\varphi_l,w) \defeq \int_0^\infty \big(r^2\varphi_l'w'+(\lambda_l+\alpha^2r^2)\varphi_l w\big)\dd r. 
	\]
We claim that $I \equiv 0$. To prove this, first note that since $\varphi_l$ solves the Euler--Lagrange equation on $(0,1)$ and $(1,\infty)$, we have
	\[
		I(\varphi_l,w) = \int_0^1 \big(r^2\varphi_l'w\big)'\dd r + \int_1^\infty \big(r^2\varphi_l'w\big)'\dd r.
	\]
Hence, for $0<\e<1<R$, we get
	\[
		\int_{\e}^R \big(r^2\varphi_l'w\big)'\dd r = R^2 \varphi_l'(R)w(R) - \e^2\varphi_l'(\e)w(\e),
	\]
since $w(1)=0$.

For $0<\e<1/2$ we estimate
	\[
		|w(\e)| \leq |w(1/2)| + \int_{\e}^{1/2} |w'|\dd r \leq |w(1/2)| + A^{1/2}(w)\left(\int_{\e}^{1/2}\frac{\dd r}{r^2}\right)^{1/2} \leq C_w \e^{-1/2}.
	\]
On the other hand, using the properties of $i_l$, $\e^2 \varphi_l'(\e)=O(\e^{l+1})$. Thus $\e^2\varphi_l'(\e)w(\e) = O(\e^{l+1/2})\to 0$ for $l\geq 1$ as $\e\to 0$.

Since $B(w)<\infty$ we have that $w\in L^2((0,\infty))$. So, there exist $R_j \in [j,2j]$ with $w^2(R_j) \leq \frac{1}{j}\int_j^{2j} w^2\dd r \to 0$. Again, a direct calculation using the properties of $k_l$ implies that $R^2\varphi_l'(R)\to 0$ as $R\to \infty$. Thus $R^2 \varphi_l'(R)w(R)\to 0$ as $R\to\infty$. Then the claim follows sending $\e\to 0$ and $R\to \infty$.

As a consequence we obtain that $Q(\varphi)=Q(\varphi_l)+Q(w)$. Since $Q(w) \geq \alpha^2\int_0^\infty w^2 r^2\dd r \geq 0$, and $Q(w)$ vanishes only when $w=0$ we conclude that $\varphi_l$ is the unique minimizer of $Q$.

Therefore, taking $\varphi_l$ instead of $w$,
	\[
		G(\lambda_l) = Q(\varphi_l) = \varphi_l'(1^-) - \varphi_l'(1^+) = \alpha \left(\frac{i_l'(\alpha)}{i_l(\alpha)}-\frac{k_l'(\alpha)}{k_l(\alpha)}\right) =\frac{1}{\alpha\,i_l(\alpha)k_l(\alpha)},
	\]
where we used the fact that 
	\[
		i_l(\alpha)k_l'(\alpha) - i_l'(\alpha)k_l(\alpha) = \frac{1}{\alpha}\big(I_{l+1/2}(\alpha)K_{l+1/2}'(\alpha) - I_{l+1/2}'(\alpha)K_{l+1/2}(\alpha)\big) = -\frac{1}{\alpha^2}
	\]
by direct computation.
\end{proof}

\end{document}
